\section[Stage 041]{Stage 041: Rank--2 Support Completion and Exact Support-Loading Theorem}
\label{stage:041}

\stagefield{Part / anchor}{Part III; primary anchor \resultanchor{MTDC-T6}.}
\claimstatus{\StatusExactClosure{} for the two-rank selected-mode determinant.}
\stagefield{Purpose}{Stage~041 adds a second loading direction for physical support and proves that the selected-branch determinant remains linear in the support loading.}
\stagefield{Inputs}{Mixed baseline \(m\), mixed direction \((1,q)\), support loading \(n\), support direction \((1,r)\), and selected lower depth \(\xi\).}

\stagefield{Verification}{SymPy audit: \StageFile{scripts/moving_throat_pde_stage041_rank2_support_sympy_audit.py}; transcript: \StageFile{scripts/output/moving_throat_pde_stage041_rank2_support_sympy_audit.txt}.  Mathematica audit: \StageFile{mathematica/moving_throat_pde_stage041_rank2_support_mathematica_audit.wl}; transcript: \StageFile{mathematica/output/moving_throat_pde_stage041_rank2_support_mathematica_audit.txt}.}

\subsection*{Derivation}

The two-direction loaded stiffness matrix is
\begin{equation}
\label{eq:app-stage041-loaded-matrix}
\frac{K_{\rm loaded}}{A_0}
=\operatorname{diag}(1,1+\delta)
-m(1,q)^T(1,q)
-n(1,r)^T(1,r).
\end{equation}
The determinant condition at \(\lambda_-=A_0(1-\xi)\) is
\begin{equation}
\label{eq:app-stage041-Dsel}
D_{\rm sel}
=\xi(\delta+\xi)
-m[\delta+(1+q^2)\xi]
-n[\delta+(1+r^2)\xi]
+mn(q-r)^2.
\end{equation}
Solving the linear equation \(D_{\rm sel}=0\) for \(n\) gives
\begin{equation}
\label{eq:app-stage041-nreq}
\boxed{
 n_{\rm req}(\xi,\delta;m,q,r)
=\frac{\xi(\delta+\xi)-m[\delta+(1+q^2)\xi]}
{\delta+(1+r^2)\xi-m(q-r)^2}}.
\end{equation}
Differentiating yields the exact monotonicity law
\begin{equation}
\label{eq:app-stage041-nreq-derivative}
\boxed{
\frac{\partial n_{\rm req}}{\partial m}
=-\frac{[\delta+(1+qr)\xi]^2}
{[\delta+(1+r^2)\xi-m(q-r)^2]^2}<0
}
\end{equation}
on every regular branch.  Additional mixed baseline always lowers the support loading needed to reach the same selected depth.

Two special closures are important:
\begin{equation}
\label{eq:app-stage041-tracking}
r=q
\quad\Longrightarrow\quad
n_{\rm req}=G_q(\xi,\delta)-m,
\end{equation}
and, with source-tied support and \(q=tR_U\),
\begin{equation}
\label{eq:app-stage041-source-tied}
 n_{\rm req}^{\rm src}
=\frac{\xi(\delta+\xi)-m[\delta+(1+\lambda_0R_U^2)\xi]}
{\delta+(1+\lambda_0)\xi-m\lambda_0(R_U-1)^2}.
\end{equation}

\stagefield{Output}{The support-loading theorem \eqref{eq:app-stage041-nreq}, its monotonicity \eqref{eq:app-stage041-nreq-derivative}, and the tracking/source-tied alternatives.}
\stagefield{Downstream use}{Stage~042 computes normalization with three distinct directions; Stage~045 shows the first coherent D/N kernel lands on the tracking surface.}

